\documentclass{article}
\usepackage[utf8]{inputenc}
\usepackage{amsmath}
\usepackage{geometry}
\geometry{margin=2cm}
\begin{document}

\section*{Reciclagem de Alumínio no Brasil em 2020}

\subsection*{Interpretação do enunciado e estratégia}
A questão apresenta dados sobre a reciclagem de latas de alumínio no Brasil em 2020, indicando que o índice de reciclagem foi de 97,4% de um total de $4{,}0 \times 10^{5}$ toneladas de latas vendidas. Sabe-se que para produzir alumínio, parte importante da matéria-prima é a bauxita, constituída de alumina (Al$_2$O$_3$). A produção industrial utiliza 50 kg de bauxita para obter 10 kg de alumínio.

O problema pede para estimar a massa de bauxita que deixou de ser extraída, devido à alta reciclagem de alumínio.

\textbf{Estratégia para resolução:}
\begin{itemize}
  \item Calcular a massa real de alumínio reciclado considerando o índice informado.
  \item Aplicar a proporção dada para encontrar a massa de bauxita correspondente, que foi poupada da extração.
\end{itemize}

\subsection*{Conteúdos e mini revisão}
\begin{table}[h!]
\centering
\begin{tabular}{|l|p{8cm}|p{6cm}|}
\hline
\textbf{Conceito} & \textbf{Ideia-chave} & \textbf{Aplicação na questão} \\
\hline
Índice de reciclagem & Percentual de latas recuperadas do total vendido & Multiplicar massa total pelas porcentagens para achar massa reciclada \\
\hline
Proporção bauxita-alumínio & 50 kg bauxita para 10 kg alumínio (5 vezes) & Multiplicar massa de alumínio reciclado por 5 para achar massa de bauxita poupada \\
\hline
\end{tabular}
\caption{Resumo dos principais conceitos envolvidos}
\end{table}

\subsection*{Resolução da Questão}
\begin{align*}
\text{Massa total de latas} &= 4,0 \times 10^{5} \, \text{ton}\newline
\text{Índice de reciclagem} &= 97{,}4\% = 0{,}974 \newline
\text{Massa de alumínio reciclado} &= 0{,}974 \times 4{,}0 \times 10^{5} = 3,896 \times 10^{5} \, \text{ton} \\
\text{Massa de bauxita poupada} &= 5 \times 3,896 \times 10^{5} = 1,948 \times 10^{6} \, \text{ton} \\
\text{Arredondando} &\to 1,9 \times 10^{6} \, \text{ton}
\end{align*}

\textbf{Pulo do gato:} Multiplicar a massa reciclada de alumínio pelo fator 5 para obter corretamente a quantidade de bauxita poupada. Aplicar o índice de reciclagem é fundamental antes dessa etapa.

\subsection*{Análise das alternativas incorretas}
\begin{itemize}
  \item \textbf{Alternativa A} subestima a quantidade por erro de escala.
  \item \textbf{Alternativa B} considera apenas alumínio reciclado sem converter em bauxita.
  \item \textbf{Alternativa C} pode refletir erro no cálculo ou arredondamento incorreto.
  \item \textbf{Alternativa E} superestima por ignorar o índice de reciclagem.
\end{itemize}

\subsection*{Código da questão}
CN\_RECICLAGEM\_TIPO\_001



\end{document}